% Source-derived chapter 08: Nearby cycles
% Original source: intersection-complexes-unramified-l-factors
\chapter{Nearby cycles}
\label{intersection-complexes-unramified-l-factors-chapter-08}
\source{intersection-complexes-unramified-l-factors}

\begin{remark}[Source section]
\label{intersection-complexes-unramified-l-factors-chapter-08-source}
\source{intersection-complexes-unramified-l-factors}
\source{intersection-complexes-unramified-l-factors}
This chapter reproduces the corresponding section of the retrieved
LaTeX source, including its definitions, statements, examples, and
proofs. It has not yet been translated into Lean.
\notready
\end{remark}

% The source section title was: Nearby cycles
\label{sect:nearby}
\source{intersection-complexes-unramified-l-factors}

\subsection{Principal degeneration} \label{sect:degeneration}
\source{intersection-complexes-unramified-l-factors}

In this section we will consider the principal degeneration $\mf X \to \mbb A^1$ 
degenerating $X$ to a horospherical variety. This is the base change of
the affine family $\mathscr X \to \ol{T_{X,\mss}}$ from \S\ref{sect:affine-degeneration}
along a certain line in the base. The principal degeneration 
was studied by \cite{Pop86} in characteristic $0$ and \cite{Grosshans92} in positive
characteristic.

We fix a choice of 
a regular dominant coweight $\ch\varrho \in \ch\Lambda^+_G \cap \ch\Lambda^\pos_G$
once and for all; none of our results depend on this choice. 
The image of $\ch\varrho$ lies in $-\Cal V$
and thus induces a map $\mbb A^1 \to \ol{T_{X,\mss}}$. 
We define the \emph{principal degeneration} 
\[ \mf X := \mathscr X \xt_{\ol{T_{X,\mss}}, \ch\varrho} \mbb A^1. \]  
We have a $G \xt \mbb G_m$-morphism 
$\mf X \to \mbb A^1$, where $G$ acts trivially on $\mbb A^1$. 
This forms an affine flat family (\cite[Proposition 9]{Pop86}). 
The fiber over $1 \in \mbb A^1$ is canonically isomorphic to 
$X$ and the fiber over $0\in \mbb A^1$ is $X_\emptyset := \Spec(\on{gr} k[X])$. 
The $\mbb G_m$-action on $\mf X$ induces an isomorphism 
$X \xt \mbb G_m \cong \mf X \xt_{\mbb A^1} \mbb G_m$. 
We also have a canonical isomorphism 
$\mf X /\!\!/ N = (X/\!\!/N) \times \mbb A^1$.

For $k$ of arbitrary characteristic, Grosshans showed that 
there is an injection 
\begin{equation}\label{e:Popov}
\source{intersection-complexes-unramified-l-factors}
    k[X_\emptyset] \to (k[X\sslash N] \ot k[N^-\bs G])^T 
\end{equation}
of $G$-algebras, which is an isomorphism 
if and only if $k[X]$ admits a $G$-module filtration with subquotients
isomorphic to dual Weyl modules of $G$ (\cite[Theorems 8, 16]{Grosshans92}). 
Of course this always holds in characteristic $0$; in positive characteristic 
we will assume that \eqref{e:Popov} is an isomorphism in what follows.

Let $\mf X^\circ$ denote the preimage of $T_X \xt \mbb A^1 \subset X/\!\!/N \xt \mbb A^1$
under the map $\mf X \to \mf X/\!\!/N$. This is the 
open subvariety which specializes 
over each fiber above $\mbb A^1$ to the dense open $B$-orbit of 
that fiber. 
Let $\mf X^\bullet \subset \mf X$ denote the open subvariety which specializes
over each fiber to the open $G$-orbit of that fiber.

\medskip

The isomorphism \eqref{e:Popov} implies that $X_\emptyset$ has a natural left $T_X$-action, 
and the orbit of the coset $N^- 1$ gives an embedding 
$T_X \into X_\emptyset$. 
We will temporarily denote its closure by $\ol{T_X}$.
\begin{lem}
\source{intersection-complexes-unramified-l-factors}
\notready  \label{lem:sectionX}
\source{intersection-complexes-unramified-l-factors}
The composition $\ol{T_X} \into X_\emptyset \to X_\emptyset/\!\!/N = X/\!\!/N$ is an 
isomorphism. In other words, we have a section
$X/\!\!/N \into X_\emptyset$. 
\end{lem}
\begin{proof}
This is essentially a special case of \cite[Proposition 7]{AT}.
For $\lambda \in \Lambda_G^+$, the
isotypic component $k[N^-\bs G]^{(\lambda)}$ is the dual Weyl module
of highest weight $\lambda$. The embedding $T \into N^-\bs G$ 
corresponds to the algebra map 
$k[N^-\bs G] \to k[T]$ that sends $k[N^-\bs G]^{(\lambda)} \to k e^\lambda$ (explicitly it sends all $T$-eigenvectors not of highest weight to zero). 
Thus, using \eqref{e:Popov}, the map $k[X_\emptyset]\to k[T_X]$ 
sends $k[X_\emptyset]^{(\lambda)} \to k e^\lambda$ for $\lambda \in \mf c_X^\vee$.
\end{proof}

\subsubsection{Contracting action} 
Let $s : X/\!\!/N\into \mf X$
denote the composition of the section given by Lemma~\ref{lem:sectionX} and the embedding
$X_\emptyset \into \mf X$ as the zero fiber.
We will construct a $\mbb G_m$-action on $\mf X$ that 
contracts $\mf X$ to the section $s$.

Recall the coweight $\ch\varrho: \mbb G_m \to T$ used to define $\mf X$. 
Let $\mbb G_m$ act on $\mf X$ via the group homomorphism 
$\mbb G_m \to G \xt \mbb G_m : a \mapsto (\ch\varrho(a^{-1}), a)$ and the
natural $G\xt \mbb G_m$-action on $\mf X$. 

\begin{lem}
\source{intersection-complexes-unramified-l-factors}
\notready \label{lem:contraction}
\source{intersection-complexes-unramified-l-factors}
The action map $\mbb G_m \times \mf X \to \mf X$ extends to a regular map 
$\mbb A^1 \times \mf X \to \mf X$ such that the composition $0 \times \mf X \to \mbb A^1 \times \mf X \to \mf X$ coincides with the composition $\mf X \to X/\!\!/N \overset{s}\to \mf X$.
\end{lem}
\begin{proof}
The action map can be described as the map of rings 
\[ k[\mf X] \to k[\mbb G_m] \ot k[\mf X] : f_\mu e^n \mapsto 
e^{n-\brac{\mu,\ch\varrho}} \ot f_\mu e^n  \]
for a $T$-eigenvector $f_\mu \in k[X]_{(\lambda)}$ of weight $\mu$. 
We have $\brac{\mu,\ch\varrho} \le \brac{\lambda,\ch\varrho} \le n$,
so the image of the co-action lies in the subalgebra 
$k[\mbb A^1] \ot k[\mf X]$. 

Moreover, since $\ch\varrho$ is regular dominant, observe that $n-\brac{\mu,\ch\varrho}=0$ 
if and only if $f_\mu$ is a highest weight vector 
and $n=\brac{\lambda,\ch\varrho}$, 
which implies that the composition $0\times \mf X \to \mbb A^1 \times \mf X \to \mf X$ 
factors through $\mf X \to X/\!\!/N \overset s \to \mf X$.
\end{proof}


\subsection{Grinberg--Kazhdan theorem in families}
The map $\mf X \to \mbb A^1$ 
induces a map of formal arc spaces $\msf L^+ \mf X \to \msf L^+ \mbb A^1$. Define
$\msf L^+_{\mbb A^1} \mf X = \msf L^+ \mf X \underset{\msf L^+ \mbb A^1}\times \mbb A^1$
where $\mbb A^1 \into \msf L^+ \mbb A^1$ is the map of constant
arcs. Then $\msf L^+_{\mbb A^1} \mf X$ is an affine scheme over $\mbb A^1$
with fiber over $\mbb G_m$ isomorphic to $\msf L^+X$ and zero fiber isomorphic to 
$\msf L^+ X_\emptyset$. 
While there does not currently exist a theory of nearby cycles for infinite type
schemes, we explain below how the nearby cycles of the ``IC complex'' of 
$\msf L^+ X$ can be modeled by nearby cycles on the global or Zastava model. 


\subsubsection{}
From now on we return to assuming that $B$ acts simply transitively on $X^\circ$, 
so $T_X=T$. Note that now \eqref{e:Popov} implies that $X_\emptyset$ is an 
affine embedding of $N^-\bs G$. 

\smallskip


First, we define the analogous models in families: let 
\[ \sM_{\mf X} = \Maps_\gen(C, \mf X/G \supset \mf X^\bullet/G), \quad 
\sY_{\mf X} = \Maps_\gen(C, \mf X/B \supset \mf X^\circ/B), \]
where $\mf X^\circ/B = \mbb A^1$.  
Since $C$ is proper, $\sM_{\mf X}$ maps to $\mbb A^1$ 
with generic fiber isomorphic to $\sM_X$ and special fiber $\sM_{X_\emptyset}$.
The arguments of \S\ref{sect:models} easily generalize to
show that $\sM_{\mf X}$ is an algebraic stack locally of finite type over $\mbb A^1$

The $B$-equivariant map $\mf X \to X\sslash N \xt \mbb A^1$ induces
a map $\sY_{\mf X} \to \sA \xt \mbb A^1$. We think of $\sY_{\mf X}\to \mbb A^1$
as a family degenerating $\sY=\sY_X$ to $\sY_\emptyset := \sY_{X_\emptyset}$. 
For $\ch\lambda\in \mf c_X$, let $\sY_{\mf X}^{\ch\lambda}$
denote the preimage of $\sA^{\ch\lambda}$ under $\pi_{\mf X} : \sY_{\mf X} \to \sA$.
We summarize the properties of $\sY_{\mf X}$ below; the proofs
are easy generalizations of those for $\sY$. 
\begin{itemize}
\item $\sY_{\mf X}^{\ch\lambda}$ is representable by a finite type scheme,
\item $\sY_{\mf X}$ satisfies the graded factorization property \emph{in families},
in the sense that there is a natural isomorphism 
\[ 
\Scr Y_{\mf X}^{\ch\lambda} \underset{\Scr A^{\ch\lambda}}\times 
(\Scr A^{\ch\lambda_1} \oo\xt \Scr A^{\ch\lambda_2}) 
\cong 
(\Scr Y_{\mf X}^{\ch\lambda_1} \underset{\mbb A^1}\times \Scr Y_{\mf X}^{\ch\lambda_2})|_{\Scr A^{\ch\lambda_1}\oo\times \Scr A^{\ch\lambda_2}} 
\]
when $\ch\lambda_1+\ch\lambda_2=\ch\lambda$.
\item There exists a closed embedding $\sY_{\mf X} \into \Gr_{B,\Sym C} \xt \mbb A^1$. 
\end{itemize}

The models $\sM_{\mf X}, \sY_{\mf X}$ are smooth-locally isomorphic 
as families, by the following generalization of Lemma~\ref{lem:localglobalyoga}: 
Let $\sY_{\mf X^\bullet} = \Maps_\gen(C, \mf X^\bullet/B \supset \mf X^\circ/B)$. 

\begin{lem}
\source{intersection-complexes-unramified-l-factors}
\notready\label{lem:yoga-fam}
\source{intersection-complexes-unramified-l-factors}
For fixed $\ch\lambda\in \mf c_X$ and any $\ch\mu\in \ch\Lambda^\pos_G$ large enough,
there is a correspondence 
\begin{equation}\label{e:Ycorrespondence}
\source{intersection-complexes-unramified-l-factors}
    \Scr Y^{\ch\lambda}_{\mf X} \leftarrow \Scr Y^{\ch\lambda}_{\mf X} \mathbin{\underset{\mbb A^1}{\oo\times}} \Scr Y^{\ch\mu}_{\mf X^\bullet} \to \Scr M_{\mf X} 
\end{equation}
over $\mbb A^1$, where the left arrow is smooth surjective and the right arrow is smooth. 
\end{lem}
The proof is the same as in \emph{loc cit.}, together with the observation that
$\Maps(C, \mf X^\bullet/G)$ is smooth because $\mf X^\bullet/G$ is the classifying stack
of a smooth group scheme over $\mbb A^1$. 

\subsubsection{}
Fix an arc $\gamma_0 \in X_\emptyset(k\tbrac t) \cap X_\emptyset^\bullet(k\lbrac t)$
and consider $\gamma_0$ as a point in $\msf L^+_{\mbb A^1} \mf X (k)$.

\begin{thm}[Grinberg--Kazhdan, Drinfeld]
\source{intersection-complexes-unramified-l-factors}
\notready  \label{thm:DGKfamily}
\source{intersection-complexes-unramified-l-factors}
There exists a point $y \in \Scr Y_\emptyset(k) \subset \Scr Y_{\mf X}(k)$
such that the formal neighborhood $(\wh{\msf L^+_{\mbb A^1} \mf X})_{\gamma_0}$ 
is isomorphic to $\wh{\mbb A}^\infty \times \wh{\Scr Y}_{\mf X, y}$. 
\end{thm}

\begin{proof}
Fix a point $v \in \abs C$ and an identification $\mf o_v \cong k\tbrac t$. 
Note that each orbit in $X_\emptyset^\bullet(F_v)/G(\mf o_v)$
has a representative in $X_\emptyset^\circ(F_v)$. 
Thus, by $G(\mf o_v)$-translation we may assume that $\gamma_0 \in X_\emptyset(\mf o_v) \cap X_\emptyset^\circ(F_v)$.
We may consider $\gamma_0$ as a section $\Spec \mf o_v \to X_\emptyset \xt^B \hat\sP_B^0$
where $\hat\sP_B^0$ is the trivial $B$-bundle on $\Spec \mf o_v$. 
By Lemma~\ref{lem:B-L}, this is equivalent to a map
$y: C \to X_\emptyset/B$ with $y(C\sm v) = \pt$. 
This is the point $y\in \sY_\emptyset(k)$ we will use. 

\smallskip

Next, we define a map of formal schemes
\begin{equation}  \label{e:LXtoY}
\source{intersection-complexes-unramified-l-factors}
(\wh{\msf L^+_{\mbb A^1} \mf X})_{\gamma_0} \to \wh{\Scr Y}_{\mf X,y} 
\end{equation}
as follows: 
Formal schemes are determined by
their $R$-points where $R$ is a local commutative
$k$-algebra with residue field $k$ whose maximal ideal $\mf m$ is nilpotent. 
Let $\gamma : \Spec R\tbrac t \to \mf X$ be an $R$-point 
of $(\wh{\msf L^+_{\mbb A^1} \mf X})_{\gamma_0}$. 
The reduction modulo $\mf m$ of $\gamma|_{\Spec R\lbrac t}$ 
equals $\gamma_0|_{\Spec k\lbrac t} \in \mf X^\circ(k\lbrac t)$. 
Since $k\lbrac t$ is the unique closed point of $R\lbrac t$, 
we deduce that $\gamma|_{\Spec R\lbrac t}$ has image contained in $\mf X^\circ$. 
Consider $\gamma$ as a section $\Spec(R \hot \mf o_v) \to \mf X \xt^B \hat\sP_B^0$
where $\hat\sP_B^0$ is the trivial $B$-bundle. 
By an easy generalization of Lemma~\ref{lem:B-L}, 
the pair $(\hat\sP^0_B, \gamma)$ is equivalent to a map 
$\tilde y : C\times \Spec(R) \to \mf X/B$
such that $\tilde y|_{(C\sm v)\times \Spec R}$
factors through $\Spec(R) \to \mbb A^1 = \mf X^\circ/B$.
We define \eqref{e:LXtoY} on $R$-points by sending 
$\gamma \mapsto \tilde y$.

Since $(\wh{\msf L^+ B})_1$ is non-canonically isomorphic to 
$\wh{\mbb A}^\infty$, the theorem follows from 
the proposition below. 
\end{proof}

\begin{prop}
\source{intersection-complexes-unramified-l-factors}
\notready\label{prop:DGKtorsor}
\source{intersection-complexes-unramified-l-factors}
Let $y : C\to X_\emptyset/B$ satisfy $y(C\sm v)=\pt$ and $y|_{\Spec \mf o_v}$ corresponds
to $(\hat\sP_B^0, \gamma_0)$. 
Then the map \eqref{e:LXtoY} is a $(\wh{\msf L^+ B})_1$-torsor. 
\end{prop}

\begin{proof}
Let $(R,\mf m)$ as above.  
By a generalization of Lemma~\ref{lem:B-L}, an $R$-point of $\wh{\Scr Y}_{\mf X,y}$ is
equivalent to a map $\bar\gamma : \Spec R \tbrac t \to \mf X/B$ such that 
$\bar\gamma|_{\Spec R\lbrac t}$ factors through $\Spec R \to \mbb A^1$. 
The fiber of \eqref{e:LXtoY} over $\bar\gamma$ parametrizes 
maps $\gamma : \Spec R \tbrac t \to \mf X$ that induce $\bar\gamma$ 
and whose reduction modulo $\mf m$ equals $\gamma_0$.
Since any $B$-bundle on $\Spec R\tbrac t$ can be trivialized, 
we see that $(\wh{\msf L^+B})_1(R)$ acts simply transitively on this fiber, since 
$\gamma_0 \in X^\circ_\emptyset(k\lbrac t)\cong B(k\lbrac t)$.
\end{proof}

\subsection{Results on nearby cycles} \label{sect:results-nearby}
\source{intersection-complexes-unramified-l-factors}
We will now consider nearby cycles on the global and Zastava models. 

\subsubsection{}
Fix $\ch\lambda \in \mathfrak c_X$ and consider the family 
$f:\Scr Y^{\ch\lambda}_{\mf X} \to \mbb A^1$. 
We have complementary embeddings 
\[ \Scr Y^{\ch\lambda}_\emptyset \overset i\into \Scr Y^{\ch\lambda}_{\mf X} \overset j\hookleftarrow \Scr Y^{\ch\lambda} \times \mbb G_m \]
where $i,j$ correspond to $f^{-1}(0), f^{-1}(\mbb G_m)$, respectively.
We let 
\[ \Psi_\sY : \rmD^b_c(\sY^{\ch\lambda} \times \mbb G_m) \to \rmD^b_c(\sY_\emptyset^{\ch\lambda})
\] 
denote the nearby cycles functor 
defined in \cite[Expos\'e XIII]{SGA7-2}, shifted by $1$ so that
$\Psi_\sY$ is perverse $t$-exact.  

Let $\Psi_\sY^u$ denote the direct summand where the monodromy operator acts unipotently. 
\begin{rem}
\source{intersection-complexes-unramified-l-factors}
\notready
Since $\mf X$ has a $\mbb G_m$-action making $f$ equivariant
and $\sF \bt \IC_{\mbb G_m}$ is $\mbb G_m$-equivariant for any sheaf $\sF$ on $\sY^{\ch\lambda}$,
a standard argument (\cite[Remark 14]{AB}) shows that $\Psi_\sY(\sF \bt \IC_{\mbb G_m}) = \Psi_\sY^u(\sF \bt \IC_{\mbb G_m})$, i.e., the monodromy is unipotent. 
\end{rem}

By $t$-exactness, $\Psi_\sY(\IC_{\sY^{\ch\lambda}\times \mbb G_m})$ is also perverse. 
In this section we will compute the image 
of $\Psi_\sY(\IC_{\sY^{\ch\lambda} \times \mbb G_m})$ in the Grothendieck group
of $\rmD^b_c(\sY^{\ch\lambda}_\emptyset)$ (equivalently, that of $\rmP(\sY^{\ch\lambda}_\emptyset)$); 
this determines the semisimplification of $\Psi_\sY(\IC_{\sY^{\ch\lambda} \times \mbb G_m})$
in the Artinian category $\rmP(\sY^{\ch\lambda}_\emptyset)$
over the algebraically closed field $k$. 

\smallskip

There is an analogous global family $\Scr M_{\mf X} \to \mbb A^1$ 
and complementary embeddings 
\[ \Scr M_{X_\emptyset} \into \Scr M_{\mf X} \hookleftarrow \Scr M_X \times \mbb G_m.  \]
Denote the associated unipotent nearby cycles functor by 
$\Psi_{\Scr M} : \rmD^b_c(\Scr M_X\times \mbb G_m) \to \rmD^b_c(\Scr M_{X_\emptyset})$.
We will simultaneously compute 
$[\Psi_{\Scr M}(\IC_{\Scr M_X\times \mbb G_m})]$.

\subsubsection{Stratifications in the horospherical case} \label{sect:horo-strat}
\source{intersection-complexes-unramified-l-factors}
Before proceeding, we give a more concrete description of the stratifications
on $\sY_\emptyset, \sM_{X_\emptyset}$ using \eqref{e:Popov}. 

Since $T_X=T$, we have $X_\emptyset^\bullet = N^-\bs G$ and 
$\Scr M_{X_\emptyset^\bullet} = \Bun_{N^-}$. 
The isomorphism \eqref{e:Popov} induces a map of affine schemes
$X/\!\!/N \times^T \ol{N^-\bs G} \to X_\emptyset$, which in turn induces
a map of stacks 
\[ \bar\iota_{\Scr M} : \Scr A \underset{\Bun_T} \times \ol\Bun_{B^-} \to \Scr M_{X_\emptyset}. \]
Let $\iota_{\Scr M} : \Scr A \underset{\Bun_T}\times \Bun_{B^-} \into \Scr M_{X_\emptyset}$ 
denote the restriction. 
For $\ch\lambda \in \mf c_X$, let $\iota^{\ch\lambda}_{\Scr M}, \bar\iota^{\ch\lambda}_{\Scr M}$ 
denote the maps corresponding to $\sA^{\ch\lambda}$.
The following is a variant of \cite[Proposition 1.2.7]{BG}, whose proof we leave to the reader.
\begin{prop}
\source{intersection-complexes-unramified-l-factors}
\notready
The map $\bar\iota_{\Scr M}^{\ch\lambda}$ is finite, and its restriction $\iota_{\Scr M}^{\ch\lambda}$ is 
a locally closed embedding. 
\end{prop}

The subschemes
$\Scr M_{X_\emptyset}^{(\ch\lambda)} = \iota_\sM^{\ch\lambda}(\Scr A^{\ch\lambda} \times_{\Bun_T^{-\ch\lambda}} \Bun_{B^-}^{-\ch\lambda})$ for $\ch\lambda\in \mf c_X$  
form a (possibly non-smooth) stratification
of $\Scr M_{X_\emptyset}$. 
This is a coarser stratification than the one defined in \S\ref{sect:globstrat}
in the following sense: 
Note that $\Cal V(X_\emptyset) = \mf t_X$ so $\mf c_{X_\emptyset}^- = \mf c_X$. 
For a partition $\mf P \in \Sym^\infty(\mf c_X\sm 0)$, 
there is a locally closed embedding $\oo C^{\mf P} \into \Scr A^{\deg(\mf P)}$, and the collection of these
subschemes ranging over all partitions $\mf P$ with $\deg(\mf P)=\ch\lambda$ forms a smooth
stratification of $\Scr A^{\ch\lambda}$. It follows from the constructions
that the strata from \S\ref{sect:globstrat} are given by 
\[ \Scr M^{\mf P}_{X_\emptyset} \cong 
\oo C^{\mf P} \underset{\Bun_T}\times \Bun_{B^-} \cong \oo C^{\mf P} \underset{\Scr A^{\deg(\mf P)}}\times \Scr M_{X_\emptyset}^{(\deg(\mf P))} \]
indexed over all $\mf P \in \Sym^\infty(\mf c_X\sm 0)$.

\medskip

Next if we consider the corresponding Zastava model, we have 
\[ \Scr Y_{X^\bullet_\emptyset} = \Scr Z^{?,0}  = \Maps(C, N^- \bs G / B \supset \pt) \]
is the \emph{open Zastava space} of Finkelberg--Mirkovi\'c. 
The \emph{Zastava space}\footnote{The Finkelberg--Mirkovi\'c Zastava space 
is the Zastava model for $\ol{G/N}$. In this paper we made a slight distinction in 
semantics between `model' and `space' to avoid confusion, but the two terms are 
interchangeable.} is in turn defined by 
$\Scr Z = \Maps_\gen(C, N^- \bs \ol{G/N} / T \supset \pt)$.
The geometry of the Zastava space has been extensively studied
in \cite{FM, FFKM, BFGM}. 
The components of $\Scr Z$ are indexed by $\ch\Lambda^\pos_G$. 

We can also define the relative open Zastava space 
${\Scr Z}^{?,0}_{\Bun_T} = \Maps_\gen(C, B^- \bs G/ B \supset \pt/T)$. 
The spaces $\Scr Z^{?,0}$ and $\Scr Z^{?,0}_{\Bun_T}$ are smooth locally isomorphic (cf.~\cite[\S 3.1]{BFGM}). 
For $\ch\lambda \in \ch\Lambda^\pos_G$, let $\Scr Z_{\Bun_T}^{\ch\lambda,0}$ denote
the preimages of $\Bun_{B^-}^{\ch\mu} \times \Bun_B^{\ch\mu-\ch\lambda}$ 
running over all $\ch\mu\in \ch\Lambda_G$. 

Since $\Scr Y_\emptyset$ is open in $\Scr M_{X_\emptyset}\times_{\Bun_G} \Bun_B$, we deduce by base change that for any $\ch\lambda\in \mf c_X, \ch\mu\in \ch\Lambda^\pos_G$, 
there is a locally closed embedding
\[ \iota_{\Scr Y_\emptyset}^{\ch\lambda,\ch\mu} : \Scr A^{\ch\lambda}\underset{\Bun_T}\times {\Scr Z}_{\Bun_T}^{\ch\mu,0} \into \Scr Y_\emptyset^{\ch\lambda+\ch\mu} \]
where we are mapping ${\Scr Z}^{?,0}_{\Bun_T} \to \Bun_{B^-} \to \Bun_T$. 
Note that ${\Scr Z}^{0,0}_{\Bun_T} = \Bun_T$, so  
$\iota^{\ch\lambda,0}_{\Scr Y_\emptyset}$ defines a map 
$\Scr A^{\ch\lambda} \into \Scr Y^{\ch\lambda}_\emptyset$. 
One can check that this map corresponds to applying $\Maps(C,T\bs ?)$ to 
the section $s_\emptyset : X/\!\!/N \into X_\emptyset$
from Lemma~\ref{lem:sectionX}. 
Therefore, 
\[ \mf s_\emptyset^{\ch\lambda} := \iota^{\ch\lambda,0}_{\Scr Y_\emptyset} : \Scr A^{\ch\lambda} \into \Scr Y^{\ch\lambda}_{\emptyset} \]
is a section of the projection $\pi_\emptyset : \Scr Y^{\ch\lambda}_\emptyset \to \Scr A^{\ch\lambda}$. 


\subsubsection{}
We compute the $*$-restriction of nearby cycles to the strata above, which
suffices to determine nearby cycles in the Grothendieck group. 
Let $\Omega(\ch\mfn_C)^{-\ch\nu} = \mbb D (\Upsilon(\ch\mfn_C)^{\ch\nu})$ 
denote the Verdier dual of the factorization algebra defined in \S\ref{sect:compareIC}. 
Recall that we defined a convolution product $\star$ on $\rmD^b_c(\sA)$ in \S\ref{sect:convolution-product}. 
The statement of our main result is: 

\begin{thm}
\source{intersection-complexes-unramified-l-factors}
\notready \label{thm:nearby}
\source{intersection-complexes-unramified-l-factors}
We have equalities in the Grothendieck group
\begin{align*}
[\Psi_\sM(\IC_{\sM_X \times \mbb G_m})] &= 
\sum_{\ch\lambda\in \mf c_X} \sum_{\ch\nu\ge 0} 
\Bigl[ \iota^{\ch\lambda}_{\sM,!}\Bigl( \bigl(\mf i_{\Scr A,\ch\nu,!}(\Omega(\ch\mfn_C)^{-\ch\nu}) \star \bar\pi_!(\IC_{\barY^{\ch\lambda-\ch\nu}}) \bigr) \bt_{\Bun_T} \IC_{\Bun_{B^-}} \Bigr) \Bigr] \\ 
[\Psi_\sY(\IC_{\sY^{\ch\mu} \times \mbb G_m})] &= 
\sum_{\substack{\ch\lambda\in \mf c_X\\ \ch\lambda \le \ch\mu}} \sum_{\ch\nu\ge 0} 
\Bigl[ \iota^{\ch\lambda,\ch\mu-\ch\lambda}_{\sY,!}\Bigl(\bigl(\mf i_{\Scr A,\ch\nu,!}(\Omega(\ch\mfn_C)^{-\ch\nu}) \star \bar\pi_!(\IC_{\barY^{\ch\lambda-\ch\nu}}) \bigr) \bt_{\Bun_T} \IC_{\Scr Z^{\ch\mu-\ch\lambda,0}_{\Bun_T}} \Bigr) \Bigr] 
\end{align*}
for any $\ch\mu\in \mf c_X$ in the second equality.
\end{thm}


We point out that the description of $\bar\pi_!\IC_{\sY^{\ch\lambda}}$ is 
the main content of the previous sections of this paper. In particular it has the format \eqref{e:format1}. 

Recall that $\Scr F\underset{\Bun_T}\boxtimes \Scr G$ denotes the $*$-restriction of 
$\Scr F\boxtimes \Scr G$ to the corresponding fiber product over $\Bun_T$,
shifted by $[-\dim \Bun_T]$. Note that $\Bun_{B^-}$ and ${\Scr Z}^{?,0}_{\Bun_T}$ are smooth stacks, so the respective IC complexes are shifted constant sheaves. 


\begin{proof}
Theorem~\ref{thm:nearby} follows by combining Corollary~\ref{cor:pibarY} 
with Lemma~\ref{lem:Psireduction} and Theorem~\ref{thm:PsiRadon} below.
\end{proof}


First, a well-known argument using some Zastava-to-global yoga 
allows us to reduce from computing restrictions to all strata to only
computing  
$\mf s^{\ch\lambda,*}_\emptyset \Psi_{\Scr Y}(\IC_{\Scr Y^{\ch\lambda}\times\mbb G_m})$
for all $\ch\lambda\in \mf c_X$. 

\begin{lem}
\source{intersection-complexes-unramified-l-factors}
\notready \label{lem:Psireduction}
\source{intersection-complexes-unramified-l-factors}
We have equalities in the Grothendieck group\footnote{In fact the isomorphisms hold in the derived category, but we omit the proof as it uses generic-Hecke equivariance to show no
twist exists on the $\Bun_{B^-}$ factor.}
\begin{align}
\label{e:YimpliesM}
\source{intersection-complexes-unramified-l-factors}
    [\iota^{\ch\lambda,*}_{\Scr M} \Psi_{\Scr M}(\IC_{\Scr M_X\times \mbb G_m})] &= [\mf s^{\ch\lambda,*}_\emptyset \bigl(\Psi_{\Scr Y}(\IC_{\Scr Y^{\ch\lambda}\times\mbb G_m})\bigr) \underset{\Bun_T}{\boxtimes} \IC_{\Bun_{B^-}}]  \\ 
\label{e:MimpliesY}
\source{intersection-complexes-unramified-l-factors}
    [\iota^{\ch\lambda,\ch\nu,*}_{\Scr Y_\emptyset} \Psi_{\Scr Y}(\IC_{\Scr Y^{\ch\lambda+\ch\nu} \times \mbb G_m})] &= [\mf s^{\ch\lambda,*}_\emptyset \bigl(\Psi_{\Scr Y}(\IC_{\Scr Y^{\ch\lambda} \times \mbb G_m}) \bigr) \underset{\Bun_T}\boxtimes \IC_{\Scr Z^{\ch\nu,0}_{\Bun_T}}],
\end{align}
where $\ch\lambda\in \mathfrak c_X, \,\ch\nu\in \ch\Lambda^\pos_G$. 
\end{lem}

\begin{proof}
The argument is the same as \cite[\S 3.1]{BFGM}, \cite[\S 8(1)]{BFGM-erratum} 
and \cite[Proof of Proposition 4.4]{BG2}. 
The strategy is that we first show \eqref{e:YimpliesM} and then use it to show \eqref{e:MimpliesY}. 

Fix $\ch\lambda'\in \mf c_X$ and $\ch\mu\in \ch\Lambda^\pos_G$ large enough as in Lemma~\ref{lem:yoga-fam} 
and consider the correspondence \eqref{e:Ycorrespondence}. 
The fiber of \eqref{e:Ycorrespondence} over $0\in \mbb A^1$
gives the correspondence
\[ \Scr Y^{\ch\lambda'}_\emptyset \leftarrow \Scr Y^{\ch\lambda'}_\emptyset \mathbin{\oo\times} {\Scr Z}^{\ch\mu,0} \to \Scr M_{X_\emptyset}. \]
Since nearby cycles commutes with smooth base change of the family over $\mbb A^1$, 
we deduce that there is an isomorphism 
\begin{equation} \label{e:PsiY=PsiM}
\source{intersection-complexes-unramified-l-factors}
 \Psi_{\Scr Y}(\IC_{\Scr Y^{\ch\lambda'}\times\mbb G_m})|^{!*}_{\Scr Y^{\ch\lambda'}_\emptyset 
\oo\times {\Scr Z}^{\ch\mu,0}} \cong \Psi_{\Scr M}(\IC_{\Scr M_X\times \mbb G_m})|^{!*}_{\Scr Y^{\ch\lambda'}_\emptyset \oo\times {\Scr Z}^{\ch\mu,0}},
\end{equation}
We now restrict to strata: for $\ch\lambda\in \mf c_X$ observe that there is a commutative diagram 
where both squares are Cartesian
\begin{equation} \label{e:cdstratacorresp}
\source{intersection-complexes-unramified-l-factors}
\begin{tikzcd}
\Scr A^{\ch\lambda} \underset{\Bun_T}\times {\Scr Z}^{\ch\lambda'-\ch\lambda,0}_{\Bun_T} 
\ar[d,hook, swap,"\iota_{\Scr Y_\emptyset}^{\ch\lambda,\ch\lambda'-\ch\lambda}"] & 
(\Scr A^{\ch\lambda} \underset{\Bun_T}\times {\Scr Z}^{\ch\lambda'-\ch\lambda,0}_{\Bun_T}) 
\mathbin{\oo\times} {\Scr Z}^{\ch\mu,0} \ar[r] \ar[d,hook] \ar[l]  &
\Scr A^{\ch\lambda} \underset{\Bun_T}\times \Bun_{B^-} \ar[d,hook, "\iota^{\ch\lambda}_{\Scr M}"] \\ 
\Scr Y^{\ch\lambda'}_\emptyset &
\Scr Y^{\ch\lambda'}_\emptyset \mathbin{\oo\times} {\Scr Z}^{\ch\mu,0} \ar[l] \ar[r] &
\Scr M_{X_\emptyset}
\end{tikzcd}
\end{equation}
By the argument of \cite[\S 8(1)]{BFGM}, every point of 
$\Scr A^{\ch\lambda} \times_{\Bun_T}\Bun_{B^-}$ is in 
the image of $\Scr A^{\ch\lambda} \mathbin{\oo\times} {\Scr Z}^{\ch\mu,0}$ 
for some $\ch\mu$ large enough, i.e., we only need to consider the diagram
\eqref{e:cdstratacorresp} when $\ch\lambda'=\ch\lambda$.
Note that $\sA^{\ch\lambda}\xt_{\Bun_T} \Bun_{B^-}$ is of finite type, so 
there exists a single $\ch\mu$ such that the map 
$\sY^{\ch\lambda}_\emptyset \oo\xt \Scr Z^{\ch\mu,0} \to \sM_{X_\emptyset}$ 
has geometrically irreducible fibers and 
contains $\Scr A^{\ch\lambda} \times_{\Bun_T}\Bun_{B^-}$ in its image (Corollary~\ref{cor:Mcover}). 
In particular, pullback along this map is fully faithful on perverse sheaves. 

By restricting the isomorphism \eqref{e:PsiY=PsiM} to the stratum
$\Scr A^{\ch\lambda} \mathbin{\oo\times} {\Scr Z}^{\ch\mu,0}$, we get 
\[
    \mf s_\emptyset^{\ch\lambda,*}\Psi_{\Scr Y}(\IC_{\Scr Y^{\ch\lambda}\times\mbb G_m})|^{!*}_{\Scr A^{\ch\lambda} \oo\times {\Scr Z}^{\ch\mu,0}} \cong \iota^{\ch\lambda,*}_{\Scr M}\Psi_{\Scr M}(\IC_{\Scr M_X\times \mbb G_m})|^{!*}_{\Scr A^{\ch\lambda} \oo\times {\Scr Z}^{\ch\mu,0}}. 
\]
This establishes the equality \eqref{e:YimpliesM} in the Grothendieck group
by fully faithfulness of the pullback.

The equality \eqref{e:MimpliesY} follows from \eqref{e:YimpliesM} and \eqref{e:PsiY=PsiM}
in the same fashion by considering the diagram \eqref{e:cdstratacorresp} with 
$\ch\lambda'=\ch\lambda+\ch\nu$.
\end{proof}

Define the functor $\Psi : \rmD^b_c(\sY^{\ch\lambda}) \to \rmD^b_c(\sY^{\ch\lambda}_\emptyset)$
by $\Psi(\sF) := \Psi_\sY(\sF \boxtimes \IC_{\mbb G_m}) = \Psi_\sY^u(\sF \bt \IC_{\mbb G_m})$. 
The crucial fact that will allow us to do our computations is the following: 

\begin{thm}
\source{intersection-complexes-unramified-l-factors}
\notready \label{thm:PsiRadon}
\source{intersection-complexes-unramified-l-factors}
There are natural isomorphisms of functors $\rmD^b_c(\sY^{\ch\lambda}) \to \rmD^b_c(\sA^{\ch\lambda})$:  
\begin{align*} 
\mf s^{\ch\lambda,*}_\emptyset \Psi 
&\cong \pi_{\emptyset,*}\Psi
\cong \pi_*  \\
\mf s^{\ch\lambda,!}_\emptyset \Psi 
&\cong \pi_{\emptyset,!}\Psi \cong \pi_!. 
\end{align*}
\end{thm}

The theorem will be proved using a standard argument
involving the contraction principle. 


\subsection{Contraction principle}

Set $\mf s^{\ch\lambda} = i \circ \mf s^{\ch\lambda}_\emptyset : \Scr A^{\ch\lambda} \into \Scr Y_{\mf X}^{\ch\lambda}$. 
We drop the superscript to denote the section $\mf s : \Scr A \into \Scr Y_{\mf X}$
on all components.
Recall that $\mf s$ corresponds to the map induced by 
the embedding $s : X/\!\!/N \to \mf X$,
and Lemma~\ref{lem:contraction} defines an action of $\mbb A^1$ on $\mf X$
that contracts to $s$. 
The action of $\mbb A^1$ commutes with that of $G$, so we get an action 
$\mbb A^1 \times \Scr Y_{\mf X} \to \Scr Y_{\mf X}$
such that $0 \times \Scr Y_{\mf X} \to\Scr Y_{\mf X}$ coincides with 
$\mf s \circ \pi_{\mf X}$. 
In this situation, the \emph{contraction principle} (\cite[Lemma 5.3]{BFGM}, 
\cite[Lemme 2.2]{VLaff}, which is closely related to Braden's theorem \cite{Braden}) says that 
there is a natural isomorphism of functors $\pi_{\mf X,*} \cong \mf s^* : \rmD^b_c(\Scr Y_{\mf X}) \to \rmD^b_c(\Scr A)$. 


\begin{proof}[Proof of Theorem \ref{thm:PsiRadon}]
We will prove the first line of isomorphisms; the second line follows from the first 
by Verdier duality. 
If we apply the contraction principle to 
$\mf s^{\ch\lambda,*}_\emptyset\Psi
= \mf s^{\ch\lambda,*} i_* \Psi$, 
we immediately get the first isomorphism 
$\mf s^{\ch\lambda,*}_\emptyset \Psi \cong \pi_{\emptyset,*}\Psi$. 


Next, we will show the isomorphism $\mf s^{\ch\lambda,*}_\emptyset \Psi \cong \pi_*$. 
Recall that \cite{Beilinson2} gives an equivalence 
$D^b \rmP(\sY^{\ch\lambda}) \cong \rmD^b_c(\sY^{\ch\lambda})$, 
so we only need to define the isomorphism on perverse sheaves. 
Let $\sF \in \rmP(\sY^{\ch\lambda})$. 
For any $a \ge 1$, let $\Scr L_a$ denote the local system on $\mbb G_m$
whose monodromy is a unipotent Jordan block of rank $a$. There are canonical
injections $\Scr L_a \to \Scr L_{a+1}$.
Beilinson's construction of the unipotent nearby cycles functor 
(cf.~\cite[2.3]{Beilinson}) gives an isomorphism 
\[ \Psi(\Scr F) \cong \colim_{a\ge 1} i^*j_*(\Scr F \bt \Scr L_a). \]
We can further apply $\mf s^{\ch\lambda,*}_\emptyset$ to get
an isomorphism $\mf s^{\ch\lambda,*}_\emptyset \Psi(\Scr F) \cong \colim \mf s^{\ch\lambda,*} j_*(\Scr F \bt \Scr L_a)$.
Applying the contraction principle, we get an isomorphism 
\[ \mf s^{\ch\lambda,*}_\emptyset \Psi(\Scr F) \cong \colim_{a\ge 1}
(\pi_{\mf X} \circ j)_*(\Scr F \bt \Scr L_a). \]
Note that $\pi_{\mf X}\circ j : \Scr Y \times \mbb G_m \to \Scr A$ 
is equal to the composition of the first projection $\Scr Y\times \mbb G_m\to \Scr Y$
and $\pi:\Scr Y \to \Scr A$. Therefore, 
\[ (\pi_{\mf X}\circ j)_* (\sF\bt \Scr L_a) 
= \pi_*(\sF) \ot H^*(\mbb G_m, \Scr L_a). \]
Since $\colim_{a\ge 1} H^*(\mbb G_m, \Scr L_a) = \ol\bbQ_\ell$, 
we conclude that $\mf s^{\ch\lambda,*}_\emptyset \Psi(\sF) \cong \pi_*(\sF)$.
\end{proof}
